\documentclass[11pt,a4paper]{article}

%% =====================================================================
%% Clay Six: Simultaneous Closure --- Phase 2 (v1, HONEST-SCOPE REWRITE)
%%
%% Title: Phase 2 --- The Simultaneous-Closure Mechanism on Canonical
%%        Encodings: Six Substrate-Level Reformulations from One
%%        Editorial Convention and a Seven-Field Bundle
%%
%% Author: Pablo Cohen (with Claude Opus 4.7 / Anthropic)
%% Date  : 2026-06-06
%% Anchor: HEAD 2b777ca / 4785c55, Lean 4 PF closure clean,
%%         zero project axioms, kernel-only
%%         [propext, Classical.choice, Quot.sound].
%% Honest-scope discipline:
%%   Every claim sourceable to a Lean file's own honest-scope. The
%%   master paper Papers/clay_substrate_reformulation_v1.tex is the
%%   model. Earlier drafts drifted into prize-claim framings; this
%%   rev corrects the drift.
%% =====================================================================

\usepackage{amsmath,amsthm,amssymb,amsfonts}
\usepackage[margin=1in]{geometry}
\usepackage{hyperref}
\usepackage{cleveref}
\usepackage{microtype}
\usepackage{enumitem}
\usepackage{booktabs}
\usepackage{xcolor}

\hypersetup{
  colorlinks=true,
  linkcolor=blue!50!black,
  urlcolor=blue!50!black,
  citecolor=blue!50!black,
}

\theoremstyle{plain}
\newtheorem{theorem}{Theorem}[section]
\newtheorem{lemma}[theorem]{Lemma}
\newtheorem{proposition}[theorem]{Proposition}
\newtheorem{corollary}[theorem]{Corollary}

\theoremstyle{definition}
\newtheorem{definition}[theorem]{Definition}
\newtheorem{solution}[theorem]{Solution}
\newtheorem{remark}[theorem]{Remark}

\title{Phase 2: The Simultaneous-Closure Mechanism on
Canonical Encodings\\
Six Substrate-Level Reformulations from One Editorial
Convention\\
and a Seven-Field Bundle\\
\large (Not a Clay Discharge of Any of the Six Axes)}
\author{Pablo Cohen\\
\small{\texttt{psolorzano@gmail.com}}\\[0.3em]
\small{Collaboration: Claude Opus 4.7 (Anthropic)}}
\date{2026-06-06}

\begin{document}
\maketitle

\begin{abstract}
This is Phase 2 of a four-paper sequence on the Principia Fractalis
(PF) framework's substrate-level reformulation of the six remaining
Clay Millennium Problems. \textbf{This is not a Clay-prize claim for
any of the six axes.} What this paper exposes is the
\emph{simultaneous-closure mechanism}: under the editorial
convention $\alpha_{\mathrm{Poincar\acute e}} := 1$ motivated by
Perelman 2003, plus a single seven-field $\mathtt{Prop}$-bundle
(\emph{five typed residuals + two $\mathtt{True}$-markers}), the six
$\mathtt{Clay\_*\_Standard}$ contracts hold simultaneously
\emph{on the framework's canonical encodings} through one Lean
theorem,
$\mathtt{perelman\_anchor\_yields\_simultaneous\_clay\_closure}$
(\texttt{PerelmanAnchoredSimultaneousClosure.lean}, line~245). Kernel
axioms: $\mathtt{[propext, Classical.choice, Quot.sound]}$; zero
project axioms.

The seven-field bundle is honestly five typed residuals plus two
$\mathtt{True}$-markers. The discharges land on the framework's
canonical encodings, which differ from the literal mathlib carriers
in named, disclosed ways: substrate-shadow vs.\ literal Chow group
(Hodge); finite-dim $\mathtt{Fin}\,2 \to \mathbb{R}$ vs.\ inf-dim
$L^{2}$ (YM); six-piece framework-$\mathtt{Prop}$ conjunction vs.\
literal Clay PDE (NS); $\mathtt{algebraicRank} := \mathtt{analyticRank}
:= \mathtt{manuscriptRankV4}$ case-split (BSD);
$\mathtt{ClassP} \neq \mathtt{ClassNP}$ on the framework's
$\mathtt{Machine}$ encoding whose two fields are free functions with
no operational Turing semantics (P vs NP); $T_{3}^{\mathrm{sym}}$ on
$\mathtt{LogWeightedL2}$ gated by Mayer 1991 \S3 and HP-program
typed residuals (RH).

The eleven framework-internal arithmetic compatibility identities on
the chosen $\alpha$-skeleton are catalogued; honest scope quoted
verbatim from the Lean file header: ``\emph{these are not Millennium
discharges. They are axiom-free algebraic facts about the 9 chosen
$\alpha$-values; the framework's interpretation is conjectural.}''
The PF\_Lean4Lean harness provides build-hash robustness across two
Lake packages (a second pass through Lean's kernel from a separate
package, \emph{not} an independent type-checker in another language).
Honest scope per the file-level honest-scope marker is preserved
throughout. Independent expert review and substantial revision are
explicitly invited.
\end{abstract}

\tableofcontents

\section{What this paper offers and what it does not}
\label{sec:offers}

This paper offers a \emph{substrate-level reformulation} of the six
remaining Clay Millennium Problems, exhibiting the
simultaneous-closure mechanism on the framework's canonical
encodings. It does \emph{not} offer a literal mathlib-tier discharge
of Riemann, P vs NP, Navier--Stokes, Yang--Mills, BSD, or Hodge at
the Clay form. The six axes remain open at that tier; the
framework's contribution is structural reformulation, not
prize-eligible discharge.

\medskip

\noindent\fbox{\parbox{0.97\textwidth}{
\textbf{What the Lean kernel carries.} A simultaneous-closure theorem
on the framework's own canonical encodings, gated by (i) the
editorial convention $\alpha_{\mathrm{Poincar\acute e}} := 1$, and
(ii) a seven-field bundle of which two fields are
$\mathtt{True}$-markers and five are typed residuals. The discharges
land on the framework's substrate-level encodings, which differ from
the literal mathlib carriers in disclosed, named ways.

\textbf{What this paper does NOT claim.} A discharge of any of the six
Clay axes at the literal mathlib tier that the Clay rules require.
The five typed residuals are unchanged.
}}

\section{The editorial convention $\alpha_{\mathrm{Poincar\acute e}} = 1$}
\label{sec:perelman-convention}

The framework adopts $\alpha_{\mathrm{Poincar\acute e}} := 1$ as its
editorial convention at the Poincar\'e axis of the chosen
$\alpha$-skeleton. This convention is \emph{motivated by} Perelman
2003's Ricci-flow proof of the Poincar\'e conjecture
\cite{perelman2002,perelman2003a,perelman2003b}. \textbf{The Lean
kernel does not consume Perelman's theorem as a hypothesis.} The
proposition $\mathtt{PerelmanAnchorInput} :
(\mathtt{framework\_alpha}).a_{\mathrm{Poincare}} = 1$ is discharged
by direct unfolding:

\begin{verbatim}
theorem perelmanAnchorInput_holds : PerelmanAnchorInput := by
  show alpha_Poincare = 1
  unfold alpha_Poincare
  norm_num
\end{verbatim}

\noindent
(\texttt{PF.Referee.PerelmanAnchoredSimultaneousClosure}, line~117).
The proof is by \texttt{norm\_num} on a definitional identity. The
Lean kernel never reads Perelman's preprints; the convention
$\alpha_{\mathrm{Poincar\acute e}} := 1$ is the framework's
identification of the substrate's Poincar\'e-axis constant, presented
under the motivation that this is the topologically-trivial axis that
Perelman 2003 made canonical at the level of three-manifold
geometrization.

\section{The eleven framework-internal arithmetic compatibility
identities}
\label{sec:invariants}

\noindent\fbox{\parbox{0.97\textwidth}{
\textbf{Honest scope (quoted verbatim from
\texttt{CrossMillenniumSharedInvariants.lean} lines 18--29):}
``\emph{These are \textbf{not} Millennium discharges. They are
axiom-free algebraic facts about the 9 chosen $\alpha$-values; the
framework's interpretation --- that each $\alpha$ is the canonical
eigenvalue parameter for its respective Millennium operator --- is
conjectural and lives elsewhere\ldots The value of this file: it
makes the \emph{interlocking} of $\alpha$-values explicit and
machine-checked, so future work cannot accidentally introduce
$\alpha$-redefinitions that break the algebraic skeleton.}''
}}

\medskip

\begin{theorem}[Eleven framework-internal arithmetic identities on
the chosen $\alpha$-skeleton]
\label{thm:eleven-invariants}
The chosen $\alpha$-skeleton constants satisfy:
\begin{align*}
&\alpha_{P}^{2} = \alpha_{\mathrm{YM}}, \quad
\alpha_{\mathrm{RH}}^{2} = 9/4, \quad
\alpha_{\mathrm{QG}}^{2} = 2\pi, \quad
\alpha_{\mathrm{Hodge}}^{2} = \alpha_{\mathrm{Hodge}} + 1, \\
&\alpha_{\mathrm{NS}} = 2\,\alpha_{\mathrm{BSD}}, \quad
\alpha_{\mathrm{NS}} = \alpha_{\mathrm{YM}}\,\alpha_{\mathrm{BSD}}, \quad
\alpha_{\mathrm{YM}} = \alpha_{\mathrm{Poincar\acute e}} + 1, \\
&\alpha_{\mathrm{RH}}\,\alpha_{\mathrm{NS}} =
  \alpha_{\mathrm{NS}} + \alpha_{\mathrm{BSD}}, \quad
\alpha_{\mathrm{RH}}\,\alpha_{\mathrm{YM}} = 3, \quad
\alpha_{NP} - \alpha_{\mathrm{Hodge}} = 1/4, \quad
\alpha_{\mathrm{QG}}^{2} = \alpha_{\mathrm{YM}}\,\pi.
\end{align*}
Each identity is provable by \texttt{ring} / \texttt{norm\_num} /
\texttt{linarith} on the definitions. Single-citation capstone:
\texttt{cross\_millennium\_shared\_invariants\_capstone}
(\texttt{PF.CrossMillenniumSharedInvariants.lean}, line~208).
\end{theorem}

\textbf{What these identities are.} Pure arithmetic compatibility
identities among the nine \emph{chosen} constants. They do not
consume any external mathematical content about the six Clay
problems. ``Forcing'' of the remaining nine $\alpha$-values from the
pinning $\alpha_{\mathrm{Poincar\acute e}} = 1$ is the uniqueness of
solutions to a definitional algebraic system, not the consumption of
external Clay content.

\begin{theorem}[Internal $\alpha$-skeleton uniqueness under the
editorial pinning]
\label{thm:alpha-forcing}
Any $\mathtt{AlphaAssignment}$ satisfying the eleven identities and
pinning $\alpha_{\mathrm{Poincar\acute e}} = 1$ equals the
framework's chosen $\alpha$-skeleton field-by-field.
\textbf{Lean:}
\texttt{PF.Referee.ClayMasterTheorem.\\
framework\_alpha\_unique\_under\_perelman\_anchor};
\texttt{PF.Referee.PerelmanAnchoredSimultaneousClosure.\\
perelman\_anchor\_yields\_alpha\_skeleton\_forcing}.
\end{theorem}

\noindent
The naming choice ``\texttt{under\_perelman\_anchor}'' is motivated
by the editorial convention. The proof object does not consume
Perelman 2003. Perelman 2003 motivates the choice of the unit value
at the Poincar\'e axis; the kernel cascade is arithmetic on the
chosen constants.

\section{The seven-field bundle: with markers flagged}
\label{sec:bundle}

The simultaneous-closure theorem consumes one bundle. The bundle has
seven fields. \textbf{Five fields are typed residuals; two fields
are $\mathtt{True}$-markers documenting symmetry, not mathematical
content.} The honest headline is therefore ``5 typed residuals + 2
symmetry markers,'' not ``7 typed residuals.''

\begin{definition}[Simultaneous Clay closure bundle, with marker fields
flagged]
\label{def:bundle}
$\mathtt{SimultaneousClayClosureBundle}$ is a seven-field
$\mathtt{Prop}$-structure. The Lean file's own comment at lines
172--175 reads: ``\emph{The two ``unconditional'' fields are
inhabited universally as \texttt{True} markers since the corresponding
Clay-Standards already hold axiom-free; they are present to make the
six-axis bundle EXPLICIT and symmetric.}'' The fields:
\begin{itemize}[leftmargin=*,nosep]
\item \textbf{(Typed residual 1, RH)} \texttt{rh\_mayer\_HP} \,:\,
  $\mathtt{Mayer1991\_SymmetricQuotientHasZetaSpectrum}$ --- Mayer
  1991 \S3 spectral-correspondence Prop.
\item \textbf{(Typed residual 2, RH)} \texttt{rh\_HP\_program} \,:\,
  $\mathtt{HilbertPolyaProgramConjecture}$ --- the published HP
  program in its standard form.
\item \textbf{(Typed residual 3, P vs NP)}
  \texttt{pnp\_classes\_distinct} \,:\, $\mathtt{ClassP} \neq
  \mathtt{ClassNP}$ on the framework's $\mathtt{Machine}$-based
  encoding (see \S\ref{sec:pnp-honest} for the encoding's
  $\mathtt{Machine}$-as-free-functions honest scope --- this field
  \emph{is} the P vs NP question on the chosen encoding).
\item \textbf{(Typed residual 4, NS)} \texttt{ns\_bootstrap} \,:\,
  $\mathtt{NS\_LocalToGlobalBootstrap}$ --- the Leray 1934 + Hopf 1951
  typed bootstrap, Fujita--Kato 1964 named residual.
\item \textbf{($\mathtt{True}$-marker 1, YM)}
  \texttt{ym\_unconditional\_marker} \,:\, $\mathtt{True}$. \\
  No mathematical content. Symmetry marker only.
\item \textbf{(Typed residual 5, BSD)}
  \texttt{bsd\_universal\_bridge} \,:\,
  $\mathtt{UniversalBridge\_MordellWeilRank\_eq\_algebraicRankV4}$
  --- the typed identity between honest mathlib
  $\mathtt{MordellWeilRank}\, E$ and the V4 framework
  $\mathtt{algebraicRankV4}\, E$ (see \S\ref{sec:bsd-honest}).
\item \textbf{($\mathtt{True}$-marker 2, Hodge)}
  \texttt{hodge\_unconditional\_marker} \,:\, $\mathtt{True}$. \\
  No mathematical content. Symmetry marker only.
\end{itemize}
\textbf{Lean:}
\texttt{PF.Referee.PerelmanAnchoredSimultaneousClosure.\\
SimultaneousClayClosureBundle} (line~176).
\end{definition}

\noindent\fbox{\parbox{0.97\textwidth}{
\textbf{Foregrounded:} the bundle field
$\mathtt{pnp\_classes\_distinct} : \mathtt{ClassP} \neq
\mathtt{ClassNP}$ \emph{is the P vs NP question itself}, on the
framework's chosen encoding. The framework's claim is that, on this
encoding, the biconditional
$\mathtt{Clay\_PvsNP\_Standard\;PF\_CanonicalComplexityEncoding} \;
\leftrightarrow\; (\mathtt{ClassP} \neq \mathtt{ClassNP})$ is provable
(line~92 of \texttt{PF.Referee.PNPCanonicalEncoding}, theorem
\texttt{Clay\_PvsNP\_Standard\_at\_canonical\_iff\_classes\_distinct}).
\emph{The biconditional is the substrate-level content.} The bundle
field is the open conjecture itself; it is not discharged in the
Lean tree.
}}

\section{The main theorem}
\label{sec:main-theorem}

\begin{theorem}[Simultaneous closure on canonical encodings]
\label{thm:main}
Under the editorial pinning $\alpha_{\mathrm{Poincar\acute e}} = 1$
and the seven-field bundle (5 typed residuals + 2
$\mathtt{True}$-markers), the six $\mathtt{Clay\_*\_Standard}$
contracts on the framework's canonical encodings hold simultaneously:
\begin{align*}
&\mathtt{Clay\_RiemannHypothesis\_Standard}, \\
&\mathtt{Clay\_PvsNP\_Standard\;PF\_CanonicalComplexityEncoding}, \\
&\mathtt{Clay\_NavierStokes\_Standard\;PF\_NS3DEncodingV4}, \\
&\mathtt{Clay\_YangMillsMassGap\_Standard\;PF\_YMEncodingV4}, \\
&\mathtt{Clay\_BSD\_Standard\;PF\_BSDEncodingV4}, \\
&\mathtt{Clay\_Hodge\_Standard\;PF\_HodgeEncoding\_FullGeneral}.
\end{align*}
\textbf{Lean:}
\texttt{PF.Referee.PerelmanAnchoredSimultaneousClosure.\\
perelman\_anchor\_yields\_simultaneous\_clay\_closure}
(line~245). Kernel axioms:
\texttt{[propext, Classical.choice, Quot.sound]}.

\medskip\noindent
\textbf{Honest scope from the file itself} (lines 44--57):
``\emph{the $\alpha$-skeleton uniqueness from Perelman's anchor is
AXIOM-FREE\ldots NOT a Clay discharge of RH or P$\neq$NP at the
literal mathlib tier. The named residuals are unchanged. The
contribution is the simultaneous-closure \emph{mechanism}.}''
\end{theorem}

\section{The six canonical encodings: per-axis honest scope}
\label{sec:six-encodings}

For each remaining Clay axis we exhibit the framework's canonical
encoding, the discharge theorem by exact Lean name, and the
honest-scope gap to the literal Clay formulation. Each subsection
quotes the Lean file's own self-description where possible.

\subsection{Riemann Hypothesis}
\label{sec:rh}

\textbf{Canonical encoding.} $T_{3}^{\mathrm{sym}}$ on
$\mathtt{LogWeightedL2}$, the framework's Mayer 1991 path.

\textbf{Discharge on the framework's encoding.}
\texttt{PF.Referee.RHCapstoneTypedBridgeV4.\\
PF\_RH\_capstone\_via\_Mayer1991\_T3sym}. Consumes
$\mathtt{rh\_mayer\_HP}$ and $\mathtt{rh\_HP\_program}$ from the
bundle.

\textbf{Honest scope.} The V4 RH bridge remains conditional on the
published Hilbert--P\'olya program in its standard form. This is not
a Clay discharge at the literal mathlib $\mathtt{riemannZeta}$ tier;
it is a discharge of the framework's $T_{3}^{\mathrm{sym}}$ encoding
conditional on the HP-program and Mayer 1991 \S3 residuals.

\subsection{P vs NP --- with the Turing-machine encoding flagged}
\label{sec:pnp}
\label{sec:pnp-honest}

\textbf{Canonical encoding.}
$\mathtt{PF\_CanonicalComplexityEncoding}$ in
$\mathtt{PF.Referee.PNPCanonicalEncoding}$ (line~81).

\textbf{Honest scope on the encoding.} The framework's
$\mathtt{Machine}$ type, defined in
$\mathtt{PrincipiaTractalis.TuringEncoding.Complexity}$ (line~71), is
\emph{not} an operational Turing machine. Its definition is verbatim:
\begin{verbatim}
structure Machine (Gamma Lambda sigma : Type) where
  accepts : BinString -> Prop
  steps   : BinString -> N
\end{verbatim}
The three type parameters $\Gamma, \Lambda, \sigma$ are
\emph{unused} in the body; the two fields $\mathtt{accepts}$ and
$\mathtt{steps}$ are \emph{free functions} with no operational
semantics constraining them to be Turing-computable. The file
docstring acknowledges this. The classes $\mathtt{ClassP}$ and
$\mathtt{ClassNP}$ built atop this structure are \emph{substrate-shadows}
of the complexity-theoretic structure, not the Cook--Karp operational
form.

\textbf{Discharge on the framework's encoding.}
\texttt{PF.Referee.PNPCanonicalEncoding.\\
Clay\_PvsNP\_Standard\_at\_canonical\_iff\_classes\_distinct}
(line~92). The unconditional content is the \emph{biconditional} of
the Clay-Standard with $\mathtt{ClassP} \neq \mathtt{ClassNP}$; the
conjecture itself is the bundle field
$\mathtt{pnp\_classes\_distinct}$.

\subsection{Navier--Stokes existence and smoothness}
\label{sec:ns}

\textbf{Canonical encoding.} $\mathtt{PF\_NS3DEncodingV4}$ in
$\mathtt{PF.NavierStokes.NS3DRegularitySolutionV4}$, with the field
$\mathtt{hasGlobalSmoothSolution := NS3DRegularitySolutionV4}$.

\textbf{Honest scope on the encoding.}
$\mathtt{NS3DRegularitySolutionV4}$ is a six-piece
framework-$\mathtt{Prop}$ conjunction (file lines 38--42): ``\emph{Six
conjuncts; the fifth conjunct is V3's BKM 1984 + Leray 1934 + Schwartz
/ forward-time + Wave 33 + Galerkin K=2 composition EXTENDED with
\texttt{LerayHopfSmoothnessConjecture} as a sixth piece.}''
\textbf{The equivalence of this six-piece conjunction to the literal
Clay Navier--Stokes statement is not established in the Lean tree.}
File HONEST SCOPE block lines 72--89: ``\emph{NOT a Clay discharge.
The literal published Leray-Hopf bootstrap is the typed encoding of
Leray 1934 + Hopf 1951 GLOBAL WEAK EXISTENCE, NOT in mathlib at HEAD
and NOT proved here\ldots the axiom-free discharge of
LerayHopfSmoothnessConjecture AT TYPED SUBSTRATE SCOPE is therefore
SOUND but does NOT touch the literal PDE Millennium content (lifting}
smoothness := True \emph{to literal} $C^{\infty}$\emph{\ldots is the
mathlib content gap)}.''

\textbf{Discharge on the framework's encoding.}
\texttt{PF.NavierStokes.NS3DRegularitySolutionV4.\\
PF\_NS\_capstone\_yields\_Clay\_NavierStokes\_standardV4}.
\emph{This is a discharge of the framework's encoding, not of the
literal Clay NS statement.}

\subsection{Yang--Mills existence and mass gap}
\label{sec:ym}

\textbf{Canonical encoding.} $\mathtt{PF\_YMEncodingV4}$ in
$\mathtt{PF.YM\_ContinuumWightmanV4}$.

\textbf{Honest scope on the encoding.} The discharge theorem
$\mathtt{PF\_YM\_capstone\_yields\_Clay\_YangMillsMassGap\_standardV4}$
(line~252) has the following file HONEST SCOPE (line~247):
``\emph{NOT a Clay discharge. Same caveats as V3 plus: the Wave 57-YM-OSRP
propagator content lives on a finite-dim} $\mathtt{Fin}\,2 \to
\mathbb{R}$ \emph{carrier, not the inf-dim} $\mathtt{L2RInf}$
\emph{carrier; the two are connected at the structural level by
shared eigenvalues} $\{1/2, 3/2\}$.'' The framework's V4 YM
Clay-Standard discharge lives on the finite-dimensional carrier; the
literal continuum SU($N$) Wightman G1--G4 axioms on an
infinite-dimensional $L^{2}$ Hilbert space are explicitly flagged as
not reached. The bundle's $\mathtt{ym\_unconditional\_marker} =
\mathtt{True}$ documents the no-residual posture
\emph{on the finite-dimensional carrier only}; the inf-dim gap is
the named open content.

\subsection{Birch--Swinnerton-Dyer}
\label{sec:bsd}
\label{sec:bsd-honest}

\textbf{Canonical encoding.} $\mathtt{PF\_BSDEncodingV4}$ in
$\mathtt{PF.Referee.BSDCapstoneTypedBridgeV4}$.

\textbf{Honest scope on the encoding.} The V4 encoding sets
$\mathtt{algebraicRank := analyticRank := manuscriptRankV4}$ (file
lines 500--505 HONEST SCOPE):
\begin{quote}
``\emph{Both rank projections share} \texttt{manuscriptRankV4}\emph{,
a case-split function surfacing the framework's discharged ranks (0,
1, 2, or 3) for 17 specific curves and} 0 \emph{for all other
curves\ldots This is NOT a Clay BSD discharge for arbitrary}
\texttt{WeierstrassCurve} $\mathbb{Q}$\emph{. For curves outside the
17-curve set, V4 returns} 0 \emph{and equality is trivially} $0 =
0$\emph{. The genuine residual is curves with rank} $\geq 1$
\emph{outside the Heegner cohort and beyond} \{E\_389a1, E\_rank\_three\},
\emph{where the framework has no axiom-free typed witness.}''
\end{quote}

The 17 specific curves carry genuine framework content. For all
other $\mathtt{WeierstrassCurve}\, \mathbb{Q}$ the encoding returns
$0$ and the equality is trivially $0 = 0$. The Clay-Standard discharge
on $\mathtt{PF\_BSDEncodingV4}$ is therefore (i)~substantive on 17
specific curves under typed published-theorem hypotheses;
(ii)~trivially $0 = 0$ on all other curves.

\textbf{Discharge on the framework's encoding.}
\texttt{PF.Referee.BSDCapstoneTypedBridgeV4.\\
PF\_BSD\_capstone\_yields\_Clay\_BSD\_standardV4} (line~506); the
body is
$\mathtt{(algebraicRankV4\_eq\_analyticRankV4\;E).symm}$ --- i.e.,
the discharge is the case-split identity. The bundle's
$\mathtt{bsd\_universal\_bridge}$ field preserves the typed identity
to honest mathlib $\mathtt{MordellWeilRank}$ where the framework has
it.

\subsection{Hodge conjecture}
\label{sec:hodge}

\textbf{Canonical encoding.} $\mathtt{PF\_HodgeEncoding\_FullGeneral}$.

\textbf{Honest scope on the encoding.} The discharge theorem
$\mathtt{pf\_hodgeEncoding\_FullGeneral\_clay\_substrate\_closure}$
works because \emph{the substrate-shadow of the Voisin obstruction is
provably False by free-coefficient matching}, not because algebraic
cycles representing transcendental Hodge classes are exhibited.
Verbatim from
\texttt{Hodge\_ClayLiteralClosureAttempt.lean} ``What this file is
NOT'' (lines 83--98):
\begin{quote}
``\emph{\textbf{NOT} a literal Clay-grade resolution of the Hodge
conjecture on general smooth projective complex varieties at codim}
$\geq 2$\emph{. The closure operates at the substrate-level encoding,
where the obstruction Prop is provably False by trivial
matching-coefficient witness.\ldots \textbf{NOT} a refutation of the
Hodge conjecture\ldots \textbf{NOT} a construction of explicit
algebraic cycles on a generic smooth quintic outside the Dwork+CM
locus (which would be a Fields-medal-grade resolution).}

\smallskip

``\emph{The gap (G1)-(G3) (lift to literal mathlib Chow +
$H^{2,2}$) is UNCHANGED.}''
\end{quote}

The bundle's $\mathtt{hodge\_unconditional\_marker} = \mathtt{True}$
documents the no-residual posture \emph{at substrate scope only}; the
codim-$\ge 2$ general-quintic content is the named open residual.

\section{The 23-problem reach record, with field-level honesty}
\label{sec:reach}

The framework develops substrate-level treatment for 23 named open
problems (7 Clay axes + 16 non-Clay open problems), bundled in
$\mathtt{framework\_universal\_reach\_realized}$
(\texttt{FrameworkUniversalReach.lean}, line~153). \textbf{At the
record level, the bundle is mostly symmetry markers.}

\noindent\fbox{\parbox{0.97\textwidth}{
\textbf{Bundle-field discipline (quoted from the Lean record's field
declarations, lines 99--137):} of the 17 fields in
$\mathtt{FrameworkUniversalReach}$, \textbf{15 are typed as}
$\mathtt{: True}$ (provenness tags carrying no mathematical content
at the record level). The two fields that \emph{do} carry non-trivial
framework Props at the record level are
$\mathtt{brocard\_attack\_realized}$ and
$\mathtt{hadwiger\_nelson\_attack\_realized}$. The substantive
per-axis content for the other 14 non-Clay problems and the seven
Clay axes lives in separate framework-attack modules
(\texttt{PF/NumberTheory/<problem>FrameworkAttack.lean}, plus the
Clay-axis-specific bridges) cited by the record file header. Reading
the record itself as a Clay-Standard discharge is incorrect; reading
it as a single citation point indexing the per-axis attack modules is
correct.
}}

\section{Independent verification posture}
\label{sec:l4l}

\subsection{Kernel-only axiom posture in the canonical tree}
\label{sec:lean-kernel}

The Lean~4 proof object for
$\mathtt{perelman\_anchor\_yields\_simultaneous\_clay\_closure}$
type-checks in the Lean~4 kernel with kernel axioms
$\mathtt{[propext, Classical.choice, Quot.sound]}$ (the standard Lean~4
kernel axioms). $\mathtt{\#print\ axioms}$ reports zero project
axioms on each capstone witness.

\subsection{The PF\_Lean4Lean harness --- what it is and what it is not}
\label{sec:l4l-honest}

\noindent\fbox{\parbox{0.97\textwidth}{
\textbf{Quoted verbatim from
\texttt{PF\_Lean4Lean/PF\_L4L/Referee/ClayVerification\\
Harness.lean}, lines~47--56:}
\begin{quote}
``\emph{This is the same Path-C protocol as
\texttt{FlagshipReverification} and
\texttt{V2AndMasterReverification}: a second pass through Lean's
kernel from a separate package with an independent build hash. It is
NOT a separate type-checker written in another language. The honest
scope of each individual closure is the one fixed by that closure's
own} \texttt{*\_honestScope} / \texttt{*\_honest\_scope}
\emph{marker in the canonical PF library (e.g.\ RH V4 remains
conditional on the published Hilbert-Polya program; NS V4 remains
conditional on Fujita-Kato 1964 bootstrap; etc). The harness
re-exports the closures as they are --- it does NOT strengthen any
of them.}''
\end{quote}
}}

\medskip

\noindent
The PF\_Lean4Lean harness provides \textbf{build-hash robustness
across two Lake packages}: the canonical theorems are re-bound in a
separate package and rebuilt from source, and the same kernel axiom
posture is re-confirmed. It does \emph{not} constitute independent
semantic verification by a separate prover; it is a second kernel
pass on the same Lean implementation. \textbf{Earlier drafts of this
paper called this an ``independent kernel re-verifier''; that framing
overstated the harness's content and is dropped here in favor of the
file's own self-description.}

\subsection{Coq mirror --- parity-only}
\label{sec:coq}

A Coq mirror provides cross-prover naming and structural parity.
Approximately 75\% of definitions in the Coq mirror are
$\mathtt{: Prop := True}$ typed mirrors of the Lean structure; the
Coq tree is \textbf{parity-only}, not independent semantic
verification of the Lean content.

\section{Empirical anchors --- honest disclosure}
\label{sec:empirical}

\subsection{The 143-problem coherence figure: retraction of
$p < 10^{-43}$ as established evidence}
\label{sec:143}

\noindent\fbox{\parbox{0.97\textwidth}{
\textbf{Retraction.} The $\mathtt{universal\_fractal\_coherence}$
Lean theorem
(\texttt{PF.Empirical.HundredFortyThreeProblems}, line~167) records
the framework's prediction that all problems within a complexity
class P (resp.\ NP) yield identical canonical $\alpha$-values. The
Lean dataset implements this as
\texttt{List.replicate 72 (canonicalEntry .P) ++
List.replicate 71 (canonicalEntry .NP)}; the
$\mathtt{coherence\_p\_value\_threshold}$ proposition is asserted,
not derived. The file's own honest-scope note (lines 36--39): ``\emph{P(all 143
by chance)} $< 10^{-43}$ \emph{(Ch 21 line 1168). We package this as
a structural} \texttt{Prop} \texttt{coherence\_p\_value\_threshold}
\emph{whose canonical witness is the manuscript's bound; it is a}
\emph{modeling assumption} \emph{about a baseline success rate, not
a derived theorem, and is labeled as such.}''

\medskip

\noindent
As currently formalized, the empirical anchor is a definitional
restatement of the framework's prediction, not a statistical analysis
of independently collected per-problem measurements. \textbf{The
$p < 10^{-43}$ figure should not be cited as established empirical
evidence}. Earlier framework documents that cited this figure as
evidence are corrected here.
}}

\subsection{IBM Quantum 9-way: framed as pre-registration}
\label{sec:ibm}

\noindent\fbox{\parbox{0.97\textwidth}{
\textbf{Pre-registration status (honest framing).} At HEAD
\texttt{4785c55} and unchanged at current HEAD, \textbf{only the pair
$(\alpha_{\mathrm{RH}}, \alpha_{NP})$ has been observed on IBM
Quantum hardware} consistent with the framework's prediction. The
other 7 $\alpha$-values (\texttt{IBMHardware9WayEvidence.lean} lines
318--320) are framework predictions \emph{committed BEFORE hardware
measurement}. The Lean theorem records the \textbf{conditional} bound
(file lines 322--327 verbatim):
\begin{quote}
``\emph{IF all 9 instances are eventually measured to within}
$10^{-3}$ \emph{of the framework's pre-specified values under the
assumed} \texttt{NoiseModel9}\emph{, THEN the random-match probability
is at most} $10^{-15}$.''
\end{quote}
\textbf{Current status:} 2 of 9 measured-consistent, 7 of 9
pre-registered, conditional bound in force. The pre-registration of
seven framework $\alpha$-values prior to hardware measurement is
scientifically valuable; it is honestly framed as such. The
conditional bound does \emph{not} say anything about the framework
absent the remaining 7 measurements.
}}

\section{Honest scope and falsifiability}
\label{sec:honest-scope}

\begin{remark}[Honest scope, foregrounded]
\label{rem:honest-scope}
This paper is \textbf{not} a literal mathlib-tier discharge of any of
the six remaining Clay axes. The substrate-level
simultaneous-closure mechanism on canonical encodings consumes the
seven-field bundle (5 typed residuals + 2 $\mathtt{True}$-markers);
the residuals are preserved openly. The six per-axis honest-scope
gaps to the literal Clay formulation are flagged in
\S\ref{sec:six-encodings}: substrate-shadow vs.\ literal Chow group
(Hodge); finite-dim $\mathtt{Fin}\,2 \to \mathbb{R}$ vs.\ inf-dim
$L^{2}$ (YM); six-piece framework-Prop conjunction vs.\ literal Clay
PDE (NS); case-split rank identity (BSD);
$\mathtt{Machine}$-as-free-functions (P vs NP); HP-program
conditional (RH). The substrate links the six structurally; this is
the framework's contribution.
\end{remark}

\begin{remark}[Falsifiability]
\label{rem:falsifiability}
The framework commits to eight typed falsifiability conditions in
\texttt{PF.Referee.FrameworkFalsifiabilityConditions}. Audit
2026-06-04: \textbf{zero of eight triggered}; \textbf{two actively
supported by 2024--2026 literature} ($F_{3}$ dark-energy direction,
$F_{6}$ $\Omega_{\Lambda} \in [0.65, 0.75]$).
\end{remark}

\section{Conclusion}
\label{sec:conclusion}

This Phase 2 paper presents the simultaneous-closure mechanism on
the framework's canonical encodings of the six remaining Clay
Millennium Problems, machine-verified in Lean~4 with kernel-only
axioms. Under the editorial convention
$\alpha_{\mathrm{Poincar\acute e}} := 1$ motivated by Perelman 2003,
plus one seven-field bundle (5 typed residuals + 2
$\mathtt{True}$-markers), the six $\mathtt{Clay\_*\_Standard}$
contracts hold simultaneously on the framework's encodings. \textbf{No
Clay-prize claim is made.} The six axes remain open at the literal
mathlib tier the Clay rules require; the framework's contribution is
structural reformulation that exhibits the six as sub-stories of one
substrate-level mechanism, with each per-axis encoding's gap to the
literal form named and preserved openly. Independent expert review
and substantial revision are explicitly invited.

The RH-axis entry vector is described in Phase 1
\cite{phase1_rh}; the P vs NP axis with its
$\mathtt{Machine}$-as-free-functions honest scope is described in
Phase 3 \cite{phase3_pnp}; the substrate-level content on
consciousness, ZPE, and cosmology is described in the Phase 4
keystone paper \cite{phase4_keystone}.

\begin{thebibliography}{99}

\bibitem{phase1_rh}
Cohen, P., Claude Opus 4.7 (2026-06-06). \emph{Phase 1 ---
A Substrate-Level Reformulation of the Riemann Hypothesis via the
Principia Fractalis Framework}.
\texttt{Papers/clay\_RH\_via\_substrate\_v2.tex}.

\bibitem{phase3_pnp}
Cohen, P., Claude Opus 4.7 (2026-06-06). \emph{Phase 3 ---
A Substrate-Level Reformulation at the P vs NP Axis via the
Principia Fractalis Framework}.
\texttt{Papers/clay\_PvsNP\_via\_substrate\_v2.tex}.

\bibitem{phase4_keystone}
Cohen, P., Claude Opus 4.7 (2026-06-06). \emph{Phase 4 Keystone ---
Substrate-Level Content on Consciousness, Zero-Point Energy, and
Cosmology via the Principia Fractalis Framework}.
\texttt{Papers/principia\_fractalis\_keystone\_v1.tex}.

\bibitem{principia_book}
Cohen, P. (2026). \emph{Principia Fractalis}, Second Edition.
HEAD of
\url{https://github.com/FractalDevTeam/Principia-Fractalis}.

\bibitem{perelman2002}
Perelman, G. (2002). The entropy formula for the Ricci flow and its
geometric applications. arXiv:math/0211159.

\bibitem{perelman2003a}
Perelman, G. (2003). Ricci flow with surgery on three-manifolds.
arXiv:math/0303109.

\bibitem{perelman2003b}
Perelman, G. (2003). Finite extinction time for the solutions to the
Ricci flow on certain three-manifolds. arXiv:math/0307245.

\bibitem{hamilton1982}
Hamilton, R. S. (1982). Three-manifolds with positive Ricci
curvature. \emph{J.\ Differential Geom.}\ 17(2): 255--306.

\bibitem{cook1971}
Cook, S.\ A. (1971). The complexity of theorem-proving procedures.
\emph{Proceedings of the Third Annual ACM Symposium on Theory of
Computing}, 151--158.

\bibitem{karp1972}
Karp, R. M. (1972). Reducibility among combinatorial problems.
In \emph{Complexity of Computer Computations}, Plenum, 85--103.

\bibitem{leray1934}
Leray, J. (1934). Sur le mouvement d'un liquide visqueux emplissant
l'espace. \emph{Acta Math.}\ 63: 193--248.

\bibitem{hopf1951}
Hopf, E. (1951). \"Uber die Anfangswertaufgabe f\"ur die
hydrodynamischen Grundgleichungen. \emph{Math.\ Nachr.}\ 4: 213--231.

\bibitem{fujitakato1964}
Fujita, H., Kato, T. (1964). On the Navier--Stokes initial value
problem. \emph{Arch.\ Rational Mech.\ Anal.}\ 16: 269--315.

\bibitem{voisin2007}
Voisin, C. (2007). Some aspects of the Hodge conjecture.
\emph{Japanese J.\ Math.}\ 2(2): 261--296.

\bibitem{mayer1991}
Mayer, D.~H. (1991). The thermodynamic formalism approach to
Selberg's zeta function for $\mathrm{PSL}(2, \mathbb{Z})$.
\emph{Bull.\ Amer.\ Math.\ Soc.}\ 25(1): 55--60.

\end{thebibliography}

\end{document}
